\section{Results and evaluation}
\label{sec:Auswertung}
%
The data taken in the experiment is written down on the measuring protocol sheets.
For convenience, the data is not transferred into tables in the results.
The protocol sheets are added to the appendix in \autoref{app:sheet} and referenced when needed.
\newline\noindent
Calculations with the resulting data are done digitally with the programming language \enquote{Python}.
Multiple packages are used, such as \enquote{Matplotlib}~\cite{matplotlib} for creating graphics, \enquote{SciPy}~\cite{scipy} for physics constants and curve-fitting and \enquote{Numpy}~\cite{numpy} for data management and math support.
%The uncertainty management as stated in \autoref{sec:fehlerrechnung} is done with the \enquote{Uncertainties}-package~\cite{uncertainties}.
%
\subsection{Standard BB84 experiment}
In the standard experiment, the probability for Alice and Bob to have chosen the same basis is $p_{\mathrm{\rotatebox[origin=c]{90}{+},\rotatebox[origin=c]{45}{+}}} = \qty{50}{\percent}$.
In every case, the ratio
\begin{equation*}
    p_{\mathrm{sifted}} = \frac{N_{\mathrm{sifted}}}{N_{\mathrm{tot}}}
\end{equation*}
between the total number of bits $N_{\mathrm{tot}}$ and the number of sifted bits $N_{\mathrm{sifted}}$ should be this probability.
\newline\noindent
The QBER from \autoref{eqn:qber} is expected be $p_{\mathrm{QBER}} = \qty{0}{\percent}$ without Eve present and $p_{\mathrm{QBER}} = \qty{25}{\percent}$ with Eve inercepting the message.
%
\subsubsection{Probabilities without Eve}
%
The bits on which the bases of Alice and Bob are the same are listed in \autoref{tab:sifted1}.
The number of sifted bits is $N_{\mathrm{sifted}} = 22$.
\begin{table}
  \centering
  \caption{Sifted bits in the standard BB84 experiment, without Eve intercepting.}
  \label{tab:sifted1}
  \sisetup{table-format=1.1, per-mode=reciprocal}
  \begin{tblr}{
      colspec = {S[table-format=2] S[table-format=1] S[table-format=1] || S[table-format=2] S[table-format=1] S[table-format=1] || S[table-format=2] S[table-format=1] S[table-format=1]},
      row{1} = {guard, mode=math},
    }
    \toprule
    \text{Bit Number} & \text{Alice} & \text{Bob} & \text{Bit Number} & \text{Alice} & \text{Bob} & \text{Bit Number} & \text{Alice} & \text{Bob} \\
    \midrule
        3  & 0 & 0 & 15 & 1 & 1 & 42 & 1 & 1 \\
        4  & 0 & 0 & 18 & 1 & 1 & 44 & 0 & 0 \\
        5  & 0 & 0 & 19 & 0 & 0 & 46 & 1 & 1 \\
        6  & 1 & 1 & 20 & 1 & 1 & 47 & 1 & 1 \\
        9  & 1 & 1 & 29 & 0 & 0 & 48 & 1 & 1 \\
        10 & 1 & 1 & 31 & 0 & 0 & 51 & 1 & 1 \\
        12 & 1 & 1 & 33 & 1 & 1 &    &   &   \\
        13 & 1 & 1 & 40 & 1 & 1 &    &   &   \\
    \bottomrule
  \end{tblr}
\end{table}
\FloatBarrier\noindent
From this, the sifted bit ratio is 
\begin{equation}
    p_{\mathrm{sifted}} = \frac{22}{52} \approx \qty{42.31}{\percent}\,.
\end{equation}
As obvious from \autoref{tab:sifted1}, the quantum bit error rate is
\begin{equation}
    p_{\mathrm{QBER}} = \qty{0}{\percent}\,,
\end{equation}
because there are no differences in the sifted bits recieved by Alice and Bob without Eve present.
\subsubsection{Probabilities with intercepting Eve}
The same data analysis is done for the experiment with Eve intercepting the bits, adding another $p = \qty{50}{\percent}$ chance of basis selection.
The number of sifted bits stays unchanged, because for Alice and Bob the bases are used.
Additionally, the same data is sent by Alice, which can be seen in comparison between the measuring protocol sheets in \autoref{app:sheet:1} and \autoref{app:sheet:3}.
The sifted bits for this part of the experiment are shown in \autoref{tab:sifted_E1}.
\begin{table}
  \centering
  \caption{Sifted bits in the standard BB84 experiment, with Eve intercepting the message.}
  \label{tab:sifted_E1}
  \sisetup{table-format=1.1, per-mode=reciprocal}
  \begin{tblr}{
      colspec = {S[table-format=2] S[table-format=1] S[table-format=1] || S[table-format=2] S[table-format=1] S[table-format=1] || S[table-format=2] S[table-format=1] S[table-format=1]},
      row{1} = {guard, mode=math},
    }
    \toprule
    \text{Bit Number} & \text{Alice} & \text{Bob} & \text{Bit Number} & \text{Alice} & \text{Bob} & \text{Bit Number} & \text{Alice} & \text{Bob} \\
    \midrule
        3  & 0 & 0 & 15 & 1 & 0 & 42 & 1 & 1 \\
        4  & 0 & 0 & 18 & 1 & 1 & 44 & 0 & 0 \\
        5  & 0 & 0 & 19 & 0 & 1 & 46 & 1 & 0 \\
        6  & 1 & 1 & 20 & 1 & 1 & 47 & 1 & 1 \\
        9  & 1 & 1 & 29 & 0 & 0 & 48 & 1 & 0 \\
        10 & 1 & 1 & 31 & 0 & 0 & 51 & 1 & 1 \\
        12 & 1 & 1 & 33 & 1 & 0 &    &   &   \\
        13 & 1 & 1 & 40 & 1 & 0 &    &   &   \\
    \bottomrule
  \end{tblr}
\end{table}
\FloatBarrier\noindent
The number of mismatching bits is $N_{\mathrm{error}} = 6$.
With \autoref{eqn:qber} this leads to a QBER of
\begin{equation}
    p_{\mathrm{QBER}} = \frac{6}{22} \approx \qty{27.27}{\percent}\,.
\end{equation}
%
\subsubsection{Transmitting a message}
With the sifted bits a key is made as a one-time pad with $N = 20$ bits length.
The word that is transmitted is \enquote{HOME}.
The encoding of the letters is done with binary addition according to the tables on \autoref{app:sheet:5}.
\newline\noindent
Using the same basis for Alice and Bob according to the sifted key, the message is transmitted from Alice to Bob then decrypted by Bob using the sifted key.
The word, Bob recieves, is indeed \enquote{HOME}.
The defining of the word, the encryption and the decryption are not written out here.
It is done in \autoref{app:sheet:1} and \autoref{app:sheet:1} and can be seen in there.
%
\subsection{BB84 with decoy states}




















